\documentclass[12pt,letterpaper]{article}
\usepackage[utf8]{inputenc}
\usepackage[spanish]{babel}
\usepackage{amsmath}
\usepackage{amsfonts}
\usepackage{amssymb}
\usepackage[margin=2.5cm]{geometry}

\begin{document}

\begin{center}
    \Large \textbf{Evaluación de Examen 2: Unidad II} \\
    \large Física para Ciencias de la Salud (005-1713) \\
    \vspace{0.5cm}
    \textbf{Estudiante:} Daniela Vallenilla \quad \textbf{C.I.:} 32.778.547
    \vspace{0.3cm} \\
    \large \textbf{Calificación Final: 9,5/10 puntos}
\end{center}

\vspace{0.5cm}

\section*{Análisis de Resultados}

\subsection*{1. Cinemática (Aceleración media)}
\textbf{Procedimiento y resultado:} Extrae apropiadamente los datos e identifica la ecuación de aceleración. Realiza la sustitución correcta restando las velocidades y dividiendo entre el tiempo ($0,6 / 0,15$). El resultado analítico es exacto: $4 \text{ m/s}^2$, acompañado de sus unidades correctas en el S.I. \\
\textbf{Veredicto:} \textbf{Correcto} (2/2 pts).

\subsection*{2. Dinámica (Leyes de Newton)}
\textbf{Procedimiento y resultado:} Muestra un despeje impecable de la aceleración partiendo de la Segunda Ley de Newton ($a = \frac{F}{m}$). Efectúa la división de $150 \text{ N} / 25 \text{ kg}$, logrando una respuesta perfectamente correcta de $6 \text{ m/s}^2$. \\
\textbf{Veredicto:} \textbf{Correcto} (2/2 pts).

\subsection*{3. Trabajo Mecánico}
\textbf{Procedimiento y resultado:} Extrae los datos y hace la excelente observación de identificar el ángulo de la fuerza respecto al movimiento ($0^\circ$). Desarrolla la sustitución obteniendo $320 \cdot 1$. No obstante, en el paso final confunde el punto de multiplicación con un signo decimal, plasmando el resultado como $320,1 \text{ Joules}$ en lugar de $320 \text{ Joules}$. \\
\textbf{Veredicto:} \textbf{Parcialmente correcto} (Se aplica una penalización mínima por el error aritmético/de notación en el último paso) (1,5/2 pts).

\subsection*{4. Energía Potencial}
\textbf{Procedimiento y resultado:} Extrae los datos (masa, altura y la constante gravitacional $9,8 \text{ m/s}^2$). Plantea el producto continuo de las tres variables según la ecuación de la energía potencial gravitatoria ($E_p = m \cdot g \cdot h$) y obtiene el resultado numérico correcto y verificado de $17,64 \text{ Joules}$. \\
\textbf{Veredicto:} \textbf{Correcto} (2/2 pts).

\subsection*{5. Potencia Mecánica}
\textbf{Procedimiento y resultado:} Plantea correctamente la relación entre trabajo mecánico y tiempo ($P = \frac{W}{t}$). Divide correctamente $75 \text{ J}$ entre $60 \text{ s}$, lo cual arroja un cálculo final preciso de $1,25 \text{ Watts}$. \\
\textbf{Veredicto:} \textbf{Correcto} (2/2 pts).

\vspace{0.5cm}
\section*{Comentarios Generales y Sugerencias}
¡Excelente examen! La estudiante evidencia un firme entendimiento de los principios físicos de la Unidad II, estructurando el desarrollo de forma clara al seguir el modelo ``Datos $\rightarrow$ Fórmula $\rightarrow$ Desarrollo''.
\begin{itemize}
    \item Se sugiere tener un especial cuidado con la **notación matemática**. Es crucial diferenciar entre un punto que indica multiplicación y uno (o coma) que denota decimales para evitar alteraciones accidentales en el resultado clínico o analítico (como sucedió en el cálculo final de la Pregunta 3).
    \item El análisis dimensional y el manejo de las fórmulas de la estudiante son sobresalientes; es muy positivo ver explícitos los subíndices y la denotación formal en cada ecuación.
    \item Se le anima a mantener el mismo empeño e impecable orden en la presentación de sus evaluaciones.
\end{itemize}

\end{document}